\section{Planificación de alto nivel guiada por lenguaje y LLM}

\subsection{Los dos papeles del LLM como planificador}

El LLM puede generar una secuencia de subtareas y también puede actuar como feasibility filter. El ponente enfatiza que el plan generado por el LLM debe llevar su provenance (hora de origen, prompt hash) y que cada step debe pasar por un tool-check.

\begin{table}[H]
\centering
\begin{tabular}{p{4cm}p{5cm}}
\toprule
Papel & Puntos prácticos \\
\midrule
Plan generator & Registrar prompt id, temperature y llm version \\
Verifier & Verificar la viabilidad del plan mediante una tool chain (por ejemplo, un python sandbox) \\
Skill mapper & Convertir el plan en skill id o primitive action \\
\bottomrule
\end{tabular}
\caption{Los múltiples papeles del LLM en un flujo jerárquico.}
\end{table}

\begin{importantbox}{LLM Prompt log}
Cada prompt lleva asociada su metadata: `timestamp`, `llm-version`, `temperature`, `tool-call`, `expected skill id`; además, la retroalimentación sobre hallucination/feasibility se devuelve al prompt log.
\end{importantbox}

\subsection{Plan verification pipeline}

\begin{figure}[H]
\centering
\includegraphics[width=0.85\textwidth]{slides-images/slide-33.jpg}
\caption{Flujo de verificación del plan del LLM y de tool-check (slide-33).}
\end{figure}

Cada salida del LLM pasa primero por el verifier (python sandbox + knowledge graph) y después el skill mapper la interpreta como un plan ejecutable. Si se detecta una hallucination, se notifica de inmediato al incident channel y el context se registra en el grade-book.

\begin{warningbox}{Guardrail de plan verification}
\begin{itemize}
    \item Cada plan incluye \texttt{plan\_id}, \texttt{llm\_version} y \texttt{prompt\_hash};
    \item Antes de la ejecución, se verifica la viabilidad con tool-check;
    \item Cuando se detecta una hallucination, se activa un rollback y se corrige el prompt.
\end{itemize}
\end{warningbox}

\subsection{Gobernanza y auditoría}

Al generar la planificación debe haber guardrails de gobernanza, como multi-round verification, tool sandbox y human-in-the-loop review. Si la LLM output contiene una high-risk action (por ejemplo, modify-files), es obligatorio activar un incident drill.

\begin{warningbox}{Riesgos de la planificación con LLM}
La planificación automatizada puede producir hallucination, sobre todo en tareas de varios pasos. Se recomienda establecer, antes del despliegue:
\begin{itemize}
    \item Multi-modal verification: verificar cada paso con una tool call (verify each step);
    \item Stop tokens: evitar que el LLM agregue subtask sin una condición de parada;
    \item Human review: confirmación humana antes de un critical plan.
\end{itemize}
\end{warningbox}

\subsection{Resumen de la sección}

El lenguaje y el LLM dan más explainability al RL jerárquico, pero también trasladan la carga de gobernanza de hallucination/incident a etapas más tempranas: el prompt log de alto nivel y el verification pipeline.
